% ========================================================
% BÀI TẬP TOÁN 5 - KẾT NỐI TRI THỨC VỚI CUỘC SỐNG
% Bài 45: Thể tích của một hình
% Định dạng: BTTeX Part III - Trắc nghiệm điền số ngắn
% ========================================================

% Câu 1
\begin{ex}
  Một hình được ghép bởi 10 hình lập phương nhỏ cạnh $1\text{ cm}$. Thể tích của hình đó là bao nhiêu xăng-ti-mét khối?
  \shortans{10}
  \loigiai{
    Thể tích mỗi hình lập phương nhỏ là 1 cm3. Thể tích hình đó là 10 cm3.
  }
\end{ex}

% Câu 2
\begin{ex}
  Một hình được ghép bởi 12 hình lập phương nhỏ cạnh $1\text{ cm}$. Thể tích của hình đó là bao nhiêu xăng-ti-mét khối?
  \shortans{12}
  \loigiai{
    Thể tích mỗi hình lập phương nhỏ là 1 cm3. Thể tích hình đó là 12 cm3.
  }
\end{ex}

% Câu 3
\begin{ex}
  Một hình được ghép bởi 14 hình lập phương nhỏ cạnh $1\text{ cm}$. Thể tích của hình đó là bao nhiêu xăng-ti-mét khối?
  \shortans{14}
  \loigiai{
    Thể tích mỗi hình lập phương nhỏ là 1 cm3. Thể tích hình đó là 14 cm3.
  }
\end{ex}

% Câu 4
\begin{ex}
  Một hình được ghép bởi 16 hình lập phương nhỏ cạnh $1\text{ cm}$. Thể tích của hình đó là bao nhiêu xăng-ti-mét khối?
  \shortans{16}
  \loigiai{
    Thể tích mỗi hình lập phương nhỏ là 1 cm3. Thể tích hình đó là 16 cm3.
  }
\end{ex}

% Câu 5
\begin{ex}
  Một hình được ghép bởi 18 hình lập phương nhỏ cạnh $1\text{ cm}$. Thể tích của hình đó là bao nhiêu xăng-ti-mét khối?
  \shortans{18}
  \loigiai{
    Thể tích mỗi hình lập phương nhỏ là 1 cm3. Thể tích hình đó là 18 cm3.
  }
\end{ex}

% Câu 6
\begin{ex}
  Một hình được ghép bởi 20 hình lập phương nhỏ cạnh $1\text{ cm}$. Thể tích của hình đó là bao nhiêu xăng-ti-mét khối?
  \shortans{20}
  \loigiai{
    Thể tích mỗi hình lập phương nhỏ là 1 cm3. Thể tích hình đó là 20 cm3.
  }
\end{ex}

% Câu 7
\begin{ex}
  Một hình được ghép bởi 22 hình lập phương nhỏ cạnh $1\text{ cm}$. Thể tích của hình đó là bao nhiêu xăng-ti-mét khối?
  \shortans{22}
  \loigiai{
    Thể tích mỗi hình lập phương nhỏ là 1 cm3. Thể tích hình đó là 22 cm3.
  }
\end{ex}

% Câu 8
\begin{ex}
  Một hình được ghép bởi 24 hình lập phương nhỏ cạnh $1\text{ cm}$. Thể tích của hình đó là bao nhiêu xăng-ti-mét khối?
  \shortans{24}
  \loigiai{
    Thể tích mỗi hình lập phương nhỏ là 1 cm3. Thể tích hình đó là 24 cm3.
  }
\end{ex}

% Câu 9
\begin{ex}
  Một hình được ghép bởi 26 hình lập phương nhỏ cạnh $1\text{ cm}$. Thể tích của hình đó là bao nhiêu xăng-ti-mét khối?
  \shortans{26}
  \loigiai{
    Thể tích mỗi hình lập phương nhỏ là 1 cm3. Thể tích hình đó là 26 cm3.
  }
\end{ex}

% Câu 10
\begin{ex}
  Một hình được ghép bởi 28 hình lập phương nhỏ cạnh $1\text{ cm}$. Thể tích của hình đó là bao nhiêu xăng-ti-mét khối?
  \shortans{28}
  \loigiai{
    Thể tích mỗi hình lập phương nhỏ là 1 cm3. Thể tích hình đó là 28 cm3.
  }
\end{ex}

% Câu 11
\begin{ex}
  Một hình được ghép bởi 30 hình lập phương nhỏ cạnh $1\text{ cm}$. Thể tích của hình đó là bao nhiêu xăng-ti-mét khối?
  \shortans{30}
  \loigiai{
    Thể tích mỗi hình lập phương nhỏ là 1 cm3. Thể tích hình đó là 30 cm3.
  }
\end{ex}

% Câu 12
\begin{ex}
  Một hình được ghép bởi 32 hình lập phương nhỏ cạnh $1\text{ cm}$. Thể tích của hình đó là bao nhiêu xăng-ti-mét khối?
  \shortans{32}
  \loigiai{
    Thể tích mỗi hình lập phương nhỏ là 1 cm3. Thể tích hình đó là 32 cm3.
  }
\end{ex}

% Câu 13
\begin{ex}
  Một hình được ghép bởi 34 hình lập phương nhỏ cạnh $1\text{ cm}$. Thể tích của hình đó là bao nhiêu xăng-ti-mét khối?
  \shortans{34}
  \loigiai{
    Thể tích mỗi hình lập phương nhỏ là 1 cm3. Thể tích hình đó là 34 cm3.
  }
\end{ex}

% Câu 14
\begin{ex}
  Một hình được ghép bởi 36 hình lập phương nhỏ cạnh $1\text{ cm}$. Thể tích của hình đó là bao nhiêu xăng-ti-mét khối?
  \shortans{36}
  \loigiai{
    Thể tích mỗi hình lập phương nhỏ là 1 cm3. Thể tích hình đó là 36 cm3.
  }
\end{ex}

% Câu 15
\begin{ex}
  Một hình được ghép bởi 38 hình lập phương nhỏ cạnh $1\text{ cm}$. Thể tích của hình đó là bao nhiêu xăng-ti-mét khối?
  \shortans{38}
  \loigiai{
    Thể tích mỗi hình lập phương nhỏ là 1 cm3. Thể tích hình đó là 38 cm3.
  }
\end{ex}

% Câu 16
\begin{ex}
  Một hình được ghép bởi 40 hình lập phương nhỏ cạnh $1\text{ cm}$. Thể tích của hình đó là bao nhiêu xăng-ti-mét khối?
  \shortans{40}
  \loigiai{
    Thể tích mỗi hình lập phương nhỏ là 1 cm3. Thể tích hình đó là 40 cm3.
  }
\end{ex}

% Câu 17
\begin{ex}
  Một hình được ghép bởi 42 hình lập phương nhỏ cạnh $1\text{ cm}$. Thể tích của hình đó là bao nhiêu xăng-ti-mét khối?
  \shortans{42}
  \loigiai{
    Thể tích mỗi hình lập phương nhỏ là 1 cm3. Thể tích hình đó là 42 cm3.
  }
\end{ex}

% Câu 18
\begin{ex}
  Một hình được ghép bởi 44 hình lập phương nhỏ cạnh $1\text{ cm}$. Thể tích của hình đó là bao nhiêu xăng-ti-mét khối?
  \shortans{44}
  \loigiai{
    Thể tích mỗi hình lập phương nhỏ là 1 cm3. Thể tích hình đó là 44 cm3.
  }
\end{ex}

% Câu 19
\begin{ex}
  Một hình được ghép bởi 46 hình lập phương nhỏ cạnh $1\text{ cm}$. Thể tích của hình đó là bao nhiêu xăng-ti-mét khối?
  \shortans{46}
  \loigiai{
    Thể tích mỗi hình lập phương nhỏ là 1 cm3. Thể tích hình đó là 46 cm3.
  }
\end{ex}

% Câu 20
\begin{ex}
  Một hình được ghép bởi 48 hình lập phương nhỏ cạnh $1\text{ cm}$. Thể tích của hình đó là bao nhiêu xăng-ti-mét khối?
  \shortans{48}
  \loigiai{
    Thể tích mỗi hình lập phương nhỏ là 1 cm3. Thể tích hình đó là 48 cm3.
  }
\end{ex}

% Câu 21
\begin{ex}
  Một hình được ghép bởi 50 hình lập phương nhỏ cạnh $1\text{ cm}$. Thể tích của hình đó là bao nhiêu xăng-ti-mét khối?
  \shortans{50}
  \loigiai{
    Thể tích mỗi hình lập phương nhỏ là 1 cm3. Thể tích hình đó là 50 cm3.
  }
\end{ex}

% Câu 22
\begin{ex}
  Một hình được ghép bởi 52 hình lập phương nhỏ cạnh $1\text{ cm}$. Thể tích của hình đó là bao nhiêu xăng-ti-mét khối?
  \shortans{52}
  \loigiai{
    Thể tích mỗi hình lập phương nhỏ là 1 cm3. Thể tích hình đó là 52 cm3.
  }
\end{ex}

% Câu 23
\begin{ex}
  Một hình được ghép bởi 54 hình lập phương nhỏ cạnh $1\text{ cm}$. Thể tích của hình đó là bao nhiêu xăng-ti-mét khối?
  \shortans{54}
  \loigiai{
    Thể tích mỗi hình lập phương nhỏ là 1 cm3. Thể tích hình đó là 54 cm3.
  }
\end{ex}

% Câu 24
\begin{ex}
  Một hình được ghép bởi 56 hình lập phương nhỏ cạnh $1\text{ cm}$. Thể tích của hình đó là bao nhiêu xăng-ti-mét khối?
  \shortans{56}
  \loigiai{
    Thể tích mỗi hình lập phương nhỏ là 1 cm3. Thể tích hình đó là 56 cm3.
  }
\end{ex}

% Câu 25
\begin{ex}
  Một hình được ghép bởi 58 hình lập phương nhỏ cạnh $1\text{ cm}$. Thể tích của hình đó là bao nhiêu xăng-ti-mét khối?
  \shortans{58}
  \loigiai{
    Thể tích mỗi hình lập phương nhỏ là 1 cm3. Thể tích hình đó là 58 cm3.
  }
\end{ex}
